\chapter{基线方法复现与 Benchmark 结果分析}

\section{FCN 骨干网络与 Source Only 基线}
本文采用一维全卷积网络作为统一特征提取 backbone。该网络由多层一维卷积、归一化与非线性映射组成，并通过时间维平均池化形成紧凑表征。Source Only 方法仅依赖源域监督分类损失进行训练，不显式处理源域与目标域之间的分布差异，因此可作为所有领域自适应方法的基础对照组。

\section{分布对齐基线方法}
\subsection{CORAL}
CORAL 通过约束源域与目标域特征协方差矩阵的一致性，减小二阶统计量差异。该方法结构简单、训练稳定，适合用来观察纯特征统计对齐在流程工业时序诊断中的迁移效果。

\subsection{DAN}
DAN 基于最大均值差异进行特征分布约束，通过核空间中的均值嵌入距离刻画源域与目标域差异。相比 CORAL，DAN 在分布对齐层面提供了更灵活的距离度量方式。

\section{对抗式基线方法}
\subsection{DANN}
DANN 在共享特征提取网络之后加入域判别器，并利用梯度反转层学习域不变表示。该方法能够直接从对抗学习角度提升跨域泛化能力，是领域自适应中最具代表性的基线之一。

\subsection{CDAN}
CDAN 在 DANN 的基础上进一步引入条件化对齐，通过结合特征表示与类别预测信息缓解简单全局对齐可能带来的类别混叠问题。该方法对于多类别故障诊断任务具有较强的参考价值。

\section{最优传输基线方法}
\subsection{DeepJDOT}
DeepJDOT 通过联合考虑特征距离与类别预测差异，建立源域与目标域样本之间的匹配关系。与仅关注边缘分布的对齐方法相比，该方法能够从样本级关系建模角度刻画跨域迁移过程。

\section{Benchmark 对比结果与问题分析}
本章实验部分将围绕统一协议下的多方法 benchmark 对比展开，重点关注目标域准确率、平衡准确率、类别混淆现象和特征空间可视化结果。当前阶段的写作重点是先固定章节结构与结果组织顺序，后续随着实验推进逐步补入主结果表格、横向比较、可视化分析和误差讨论。

\section{本章小结}
本章围绕 FCN 统一 backbone 下的各类领域自适应基线方法，给出了 benchmark 复现章节的组织框架。后续实验结果将按照 Source Only、分布对齐、对抗式方法和最优传输方法的顺序展开，并结合总体结果与可视化分析讨论多工况迁移中的主要难点。
